\documentclass[10pt]{article}

% ---------- Document Identity (update these only) ----------
\newcommand{\DocStage}{}
\newcommand{\DocTitle}{Your Road to Tier One SOF}
\newcommand{\ProgramName}{182-Week Civilian-to-Selection Readiness Pipeline}
\newcommand{\ProgramSubtitle}{}

% ---------- Page ----------
\usepackage[legalpaper,left=0.5in,right=0.5in,top=0.6in,bottom=0.45in]{geometry}

% ---------- Typography ----------
\usepackage[T1]{fontenc}
\usepackage[utf8]{inputenc}
\usepackage{libertinus}
\usepackage{microtype}
\usepackage{parskip}
\usepackage{ragged2e}
\usepackage{textcomp}
\AtBeginDocument{\color{black}}
\AtBeginDocument{\pagecolor{white}}
\AtBeginDocument{\justifying}

% ---------- Landscape pages ----------
\usepackage{pdflscape}

% ---------- Color ----------
\usepackage[table]{xcolor}
\definecolor{StageBlue}{HTML}{1F4E79}
\definecolor{StageBlueDark}{HTML}{163A5A}
\definecolor{LightGray}{HTML}{F2F4F7}
\definecolor{MidGray}{HTML}{D9DEE5}
\definecolor{RuleGray}{HTML}{C7CED8}
\definecolor{HeaderInk}{HTML}{000000}
\definecolor{Ink}{HTML}{000000}

% ---------- Utility ----------
\usepackage{enumitem}
\sloppy
\setlength{\emergencystretch}{2em}

% ---------- Boxes / UI ----------
\usepackage[most]{tcolorbox}

\tcbset{
  charterTitle/.style={
    enhanced,
    colback=StageBlue!4,
    colframe=StageBlue,
    boxrule=1.25pt,
    arc=3pt,
    left=16pt,right=16pt,top=14pt,bottom=14pt,
    borderline north={3.2pt}{0pt}{StageBlue},
    borderline west={4pt}{0pt}{StageBlue},
    borderline south={1.25pt}{0pt}{StageBlue},
    drop shadow={black!25!white},
  },
  charterRule/.style={
    enhanced,
    colback=LightGray,
    colframe=RuleGray,
    boxrule=0.85pt,
    arc=3pt,
    left=12pt,right=12pt,top=10pt,bottom=10pt,
    borderline west={3pt}{0pt}{StageBlue},
  },
  charterBadge/.style={
    enhanced,
    colback=StageBlue,
    coltext=white,
    colframe=StageBlueDark,
    boxrule=0.6pt,
    arc=2pt,
    left=6pt,right=6pt,top=3pt,bottom=3pt,
  }
}
\newtcbox{\Badge}{nobeforeafter,charterBadge}

% ---------- Lists ----------
\setlist[itemize]{leftmargin=1.5em, itemsep=0.25em, topsep=0.25em}
\setlist[enumerate]{leftmargin=1.8em, itemsep=0.25em, topsep=0.25em}

% ---------- TOC ----------
\usepackage{tocloft}
\setcounter{tocdepth}{3}
\setlength{\cftbeforesecskip}{3pt}
\setlength{\cftbeforesubsecskip}{2pt}
\setlength{\cftbeforesubsubsecskip}{0pt}
\renewcommand{\cftsecfont}{\bfseries}
\renewcommand{\cftsecpagefont}{\bfseries}
\renewcommand{\cftsubsecfont}{\normalfont}
\renewcommand{\cftsubsecpagefont}{\normalfont}
\renewcommand{\cftsubsubsecfont}{\normalfont}
\renewcommand{\cftsubsubsecpagefont}{\normalfont}
\setlength{\cftsecindent}{0pt}
\setlength{\cftsubsecindent}{1.25em}
\setlength{\cftsubsubsecindent}{2.8em}
\setlength{\cftsecnumwidth}{2.1em}
\setlength{\cftsubsecnumwidth}{2.8em}
\setlength{\cftsubsubsecnumwidth}{3.6em}

% Remove the built-in TOC heading so only the custom heading is shown
\makeatletter
\renewcommand{\tableofcontents}{\@starttoc{toc}}
\makeatother

% ---------- Header/Footer: page number only ----------
\usepackage{fancyhdr}
\pagestyle{fancy}
\fancyhf{}
\renewcommand{\headrulewidth}{0pt}
\renewcommand{\footrulewidth}{0pt}
\fancyfoot[C]{\thepage}

% ---------- Section formatting ----------
\usepackage{titlesec}

\titleformat{\section}
  {\Large\bfseries\color{Ink}}
  {\thesection}{0.6em}{}
  [\vspace{0.25em}\textcolor{StageBlue}{\rule{\linewidth}{0.8pt}}]

\titleformat{\subsection}
  {\large\bfseries\color{Ink}}
  {\thesubsection}{0.6em}{}

\titleformat{\subsubsection}
  {\normalsize\bfseries\color{Ink}}
  {\thesubsubsection}{0.6em}{}

\titlespacing*{\section}{0pt}{0.95em}{0.35em}
\titlespacing*{\subsection}{0pt}{0.65em}{0.25em}
\titlespacing*{\subsubsection}{0pt}{0.55em}{0.2em}

% ---------- Table engine ----------
\usepackage{tabularray}
\UseTblrLibrary{booktabs}

% ---------- tabularray: GLOBAL table treatment ----------
\SetTblrInner{
  cells = {fg=black, bg=white, valign=t},
}

% ---------- Header helpers ----------
\newcommand{\Hdr}[1]{\shortstack[c]{\strut #1\strut}}

% ---------- Lines ----------
\newcommand{\TitleLine}{\textcolor{StageBlue}{\rule{0.98\linewidth}{0.95pt}}}
\newcommand{\SubTitleLine}{\textcolor{RuleGray}{\rule{0.98\linewidth}{0.65pt}}}

% ---------- Continued on next page ----------
\newcommand{\ContinuedOnNextPageInline}{%
  \par
  \vfill
  \noindent\textcolor{RuleGray}{\rule{\linewidth}{1pt}}\vspace{-2pt}
  \noindent\hfill\textit{Continued on next page}\par\vspace{-10pt}
  \noindent\textcolor{RuleGray}{\rule{\linewidth}{1pt}}
  \clearpage
}

% ---------- PDF hyperlinks ----------
\usepackage[hidelinks]{hyperref}
\hypersetup{
  colorlinks=false,
  pdfborder={0 0 0},
  linktoc=all,
  pdftitle={\DocTitle},
  pdfauthor={\ProgramName}
}

% =========================================================
\begin{document}

\begin{landscape}

\subsection*{6.B.1 Matrix B1 — Hardware / Environment / Systems}
\phantomsection
\addcontentsline{toc}{subsection}{6.B.1 Matrix B1 — Hardware / Environment / Systems}

Legend: T0/T1/T2/T3 = Tier 0–3 minimum capability by phase end gate. Timing tokens: (E) required by phase entry, (M) required by mid-phase gate window, (End) required by end-phase gate window. If no token is shown, assume required by end-phase gate window. Stage 3 protected Deload+Test windows anchor timing interpretation; “no novelty near Deload+Test” applies to any first-introduction.

\begingroup
\scriptsize
\setlength{\tabcolsep}{2pt}
\renewcommand{\arraystretch}{1.12}

\begin{longtblr}{
  width=\linewidth,
  rowhead=1,
  colspec = {
    Q[l,wd=0.40in]
    Q[l,wd=1.0in]
    Q[l,wd=0.35in]
    Q[l,wd=0.35in]
    Q[l,wd=0.35in]
    Q[l,wd=0.35in]
    Q[l,wd=0.35in]
    Q[l,wd=0.35in]
    Q[l,wd=0.35in]
    Q[l,wd=0.35in]
    Q[l,wd=0.35in]
    Q[l,wd=0.35in]
    Q[l,wd=0.35in]
    Q[l,wd=0.35in]
    Q[l,wd=0.35in]
    Q[l,wd=0.35in]
    Q[l,wd=0.35in]
    X[l,1.00]
  },
  hlines,
  vlines,
  cells = {hsep=6pt, vsep=6pt, halign=l, valign=t},
  row{odd} = {bg=white},
  row{even} = {bg=LightGray},
  row{1} = {bg=StageBlue!15, fg=black, font=\bfseries, halign=l, valign=m},
}
\Hdr{Phase #} &
\Hdr{Phase Name} &
\Hdr{7.01} &
\Hdr{7.08} &
\Hdr{7.07} &
\Hdr{7.11} &
\Hdr{7.12} &
\Hdr{7.09} &
\Hdr{7.10} &
\Hdr{7.13} &
\Hdr{7.14} &
\Hdr{7.15} &
\Hdr{7.16} &
\Hdr{7.17} &
\Hdr{7.18} &
\Hdr{7.19} &
\Hdr{7.20} &
\Hdr{Notes New Additions / Milestones} \\

00 &
Bodily Reactivation and Baseline Creation &
T0 &
T0 &
T0 &
T1 (E) &
T1 (E) &
T0 &
T1 (End) &
T1 (E) &
T1 (E) &
T1 (E) &
T1 (E) &
T0 &
T1 (E) &
T0 &
T1 (E) &
Phase 00 minimal posture enforced: no heavy load carriage requirement; no underwater; no maximal strength equipment requirement. First hard dependencies for execution/admissibility: 7.11 footwear interface and footcare basics; 7.18 timing/measurement/logging; 7.13 hydration container access; basic 7.12 safe surface + stable chair/bench for low-risk movement; 7.14/7.15 apparel baseline to prevent environment-driven noncompliance; 7.20 hygiene/drying posture. 7.10 logistics reaches Tier 1 by entry to prevent repeated wet-footwear/skin breakdown and storage friction. \\

01 &
Base Conditioning and Hypertrophy &
T1 (End) &
T1 (End) &
T0 &
T1 &
T1 &
T0 &
T1 &
T1 &
T1 &
T1 &
T1 &
T0 &
T1 &
T0 &
T1 &
First minimum deployable resistance capability (7.01 Tier 1) becomes required by end gate to support base strength/hypertrophy intent without gym dependency (generic resistance capability; not barbell/rack-level). 7.08 Tier 1 becomes required by end gate for year-round low-impact continuity and indoor substitution when outdoor conditions compromise safety/validity. Testing/timing (7.18) remains Tier 1 to support repeated \textbf{CAL}/\textbf{RUN} gate days under controlled conditions and “no novelty” posture. \\

02 &
Maximal Strength Development &
T1 (E) &
T1 &
T0 &
T1 &
T1 &
T0 &
T1 &
T1 &
T1 &
T1 &
T1 &
T0 &
T1 &
T0 &
T1 &
First barbell/rack-level strength capability may become functionally Tier 1 in Phase 02 only when explicitly required and fully reconciled across Stage 6 and Stage 7. Maintain strict “no novelty near Deload+Test” for any strength implement changes that would affect comparability and safety. \\

03 &
Converting Strength to Power and Endurance for SOF Selection &
T1 &
T1 &
T1 (End) &
T1 &
T1 &
T0 &
T1 &
T1 &
T1 &
T1 &
T1 &
T0 &
T2 (End) &
T1 (End) &
T1 &
Locked milestone: 7.07 reaches Tier 1 first in Phase 03 by end gate (minimum deployable load-bearing/carry capability established for the next phase’s load-tolerance demands; configuration must be loggable and stable). 7.18 upgrades to Tier 2 by end gate to support more complex, repeatable gate days (route verification posture, load verification posture, and tighter tool standardization). 7.19 reaches Tier 1 by end gate where protective/assistive capability is needed to keep durability stable as complexity and consequences rise. \\

04 &
Strength-Endurance and Work Capacity Development &
T1 &
T1 &
T1 (E) &
T1 &
T1 &
T0 &
T1 &
T1 &
T1 &
T1 &
T1 &
T1 (End) &
T2 &
T1 &
T1 &
Locked milestone: 7.17 first becomes Tier 1 in Phase 04 (surface-only aquatic capability readiness and safety logistics become deployable by end gate; underwater remains governed and not implied). Ruck capability remains Tier 1 and is treated as stable configuration input; any configuration changes near protected windows are treated as novelty and require stabilization/reslot. \\

05 &
Aerobic Base and Extended Stamina Development &
T1 &
T1 &
T1 &
T1 &
T1 &
T0 &
T1 &
T1 &
T1 &
T1 &
T1 &
T1 &
T2 &
T1 &
T1 &
First phase where extended endurance verification is typically consequential; logistics and environment continuity become primary risk controls (7.10/7.14 remain Tier 1 minimum). Testing/timing remains Tier 2 to preserve comparability across longer events and to prevent route/tool drift from contaminating gate evidence. \\

06 &
Lactate Threshold and Power-Endurance Development &
T2 (End) &
T2 (End) &
T2 (End) &
T2 (End) &
T2 (End) &
T1 (End) &
T2 (End) &
T2 (End) &
T2 (End) &
T2 (End) &
T1 &
T2 (End) &
T2 (End) &
T2 &
T2 (End) &
First major “environment volatility becomes constant” point for Grand Forks execution continuity: 7.14 Tier 2 by end gate (expanded environmental coverage) is the first climate-driven milestone where inadequate apparel becomes a repeated validity/safety failure risk rather than a minor inconvenience. 7.09 becomes Tier 1 by end gate as minimum deployable monitoring/log visibility support for high-consequence load management (non-diagnostic; trend support only). \\

07 &
Advanced Hybrid Overload (U.S. SOF Transformation Program) &
T2 &
T2 &
T2 &
T2 &
T2 &
T1 &
T2 &
T2 &
T2 &
T2 &
T1 &
T2 &
T2 &
T2 &
T2 &
Hybrid density increases the cost of friction: Tier 2 sustainment posture (7.10) and Tier 2 self-management (7.16) become functionally mandatory to prevent schedule drift, undocumented substitutions, and test-week validity contamination. Ruck and feet remain Tier 2 due to compounding foot interface and load tolerance demands. \\

08 &
Structural Durability and Load Tolerance &
T2 &
T2 &
T2 &
T2 &
T2 &
T1 &
T2 &
T2 &
T2 &
T2 &
T1 &
T2 &
T2 &
T2 &
T2 &
Load tolerance phases treat 7.07/7.11/7.19 Tier 2 as a stability requirement: fit/adjustment repeatability, footcare redundancy, and protective posture reduce predictable breakdown. No “upgrade-by-novelty” near protected windows; changes require stabilization or reslot. \\

09 &
High-Volume Calisthenics and Stamina &
T2 &
T2 &
T2 &
T2 &
T2 &
T1 &
T2 &
T2 &
T2 &
T2 &
T1 &
T2 &
T2 &
T2 &
T2 &
High \textbf{CAL} volume increases skin/hand/foot friction risks and recovery needs; Tier 2 recovery/mobility tooling remains minimum deployable. Testing/timing remains Tier 2 to keep artifact integrity and comparability stable through increased density. \\

10 &
Advanced Tactical Readiness &
T2 &
T2 &
T2 &
T2 &
T2 &
T1 &
T2 &
T2 &
T2 &
T2 &
T1 &
T2 &
T2 &
T2 &
T2 &
Multi-domain readiness demands stable environment, course-setup, and sustainment posture; underwater remains conditional and governance-gated (not implied by Tier level). Aquatic category remains Tier 2 as a surface-and-safety capability baseline; underwater prerequisites are handled by governance and do not become an equipment-driven requirement. \\

11 &
Pre-Selection Conditioning and Recovery &
T2 &
T2 &
T2 &
T2 &
T2 &
T1 &
T2 &
T2 &
T2 &
T2 &
T1 &
T2 &
T2 &
T2 &
T2 &
Recovery sensitivity rises; equipment posture must reduce friction and prevent last-minute substitutions. Tier 2 sustainment and recovery tooling remain minimum deployable to avoid avoidable sleep/foot/skin breakdown that invalidates final verification windows. \\

12 &
Final Pre-Selection Conditioning &
T2 &
T2 &
T3 (End) &
T3 (End) &
T3 (End) &
T1 &
T3 (End) &
T2 &
T2 &
T2 &
T1 &
T3 (End) &
T3 (End) &
T3 (End) &
T3 (End) &
First “overbuild for verification reliability” milestone: Tier 3 is introduced selectively by end gate where redundancy and setup repeatability protect final evidence quality. Tier 3 emphasis is on reducing variance (backup measurement capability, stable course-setup, redundant ruck/foot interface options, and recovery tooling redundancy), not on purchasing for performance. \\

13 &
Full-Spectrum Readiness Consolidation &
T2 &
T2 &
T3 &
T3 &
T3 &
T1 &
T3 &
T2 &
T2 &
T2 &
T1 &
T3 &
T3 &
T3 &
T3 &
Consolidation phase maintains Tier 3 verification reliability posture established in Phase 12 for the categories that directly affect repeatability and admissibility of final gates (7.18, 7.07, 7.11, 7.12, 7.10, 7.16, 7.19, 7.17). Any late-stage capability change near protected windows is treated as novelty and triggers reslot rather than “quick adaptation.” \\

\end{longtblr}

\endgroup

\clearpage

\subsection*{6.B.2 Matrix B2 — Supplements / Support}
\phantomsection
\addcontentsline{toc}{subsection}{6.B.2 Matrix B2 — Supplements / Support}

\begingroup
\scriptsize
\setlength{\tabcolsep}{2pt}
\renewcommand{\arraystretch}{1.12}

\begin{longtblr}{
  width=\linewidth,
  rowhead=1,
  colspec = {
    Q[l,wd=0.3in]
    Q[l,wd=2.35in]
    Q[l,wd=0.30in]
    Q[l,wd=0.30in]
    Q[l,wd=0.30in]
    Q[l,wd=0.30in]
    Q[l,wd=0.95in]
    X[l,3.45]
  },
  hlines,
  vlines,
  cells = {hsep=6pt, vsep=6pt, halign=l, valign=t},
  row{odd} = {bg=white},
  row{even} = {bg=LightGray},
  row{1} = {bg=StageBlue!15, fg=black, font=\bfseries, halign=l, valign=m},
}
\Hdr{Phase \#} &
\Hdr{Phase Name} &
\Hdr{7.03} &
\Hdr{7.06} &
\Hdr{7.05} &
\Hdr{7.04} &
\Hdr{7.02} &
\Hdr{Notes Deployment Logic / Stability Rules} \\

00 &
Bodily Reactivation and Baseline Creation &
T1 (E) &
T1 (E) &
T0 &
T0 &
T0 &
Baseline stack posture becomes legitimate here: simple, stable foundational health support and simple recovery/sleep support may be deployed from entry when individually tolerated and logged consistently. No stimulant posture and no performance-aid escalation while tolerance, sleep stability, GI stability, and adherence are still being established. \\

01 &
Base Conditioning and Hypertrophy &
T1 &
T1 &
T0 &
T1 (End) &
T0 &
First phase where 7.04 Tier 1 becomes justified: base conditioning and hypertrophy can support a stable, low-complexity performance-support posture if it is tolerable and does not create sleep drift, GI disturbance, appetite disruption, or weigh-in noise. Foundational and recovery support remain the primary minimum. \\

02 &
Maximal Strength Development &
T1 &
T1 &
T0 &
T1 &
T0 &
Strength phase rewards stability rather than novelty. Foundational support, recovery support, and a simple performance-support baseline may remain in place, but the stack must not change near proxy or gate windows in ways that alter comparability, hydration behavior, or next-day readiness. \\

03 &
Converting Strength to Power and Endurance for SOF Selection &
T1 &
T1 &
T1 (End) &
T1 &
T0 &
Controlled energy-enhancer posture first becomes legitimate by end gate: the issue is not larger stimulant exposure, but standardized use, cutoff discipline, and no surprise variability on test weeks. Performance support remains subordinate to recovery and logging discipline. \\

04 &
Strength-Endurance and Work Capacity Development &
T1 &
T1 &
T1 &
T1 &
T0 &
Work-capacity phase makes a stable pre-session energy posture operationally relevant, but only if it is repeatable and does not impair sleep, hydration, appetite control, or water-governance compliance. Support stack must remain simple enough that missed doses or substitutions do not destabilize protected windows. \\

05 &
Aerobic Base and Extended Stamina Development &
T1 &
T1 &
T1 &
T1 &
T0 &
Longer aerobic work increases the value of stable foundational, recovery, performance, and controlled energy support. Hydration and fueling discipline still dominate; supplements support the plan and do not replace it. Any first-introduction near Deload+Test remains a novelty event and is not admissible as “part of the plan.” \\

06 &
Lactate Threshold and Power-Endurance Development &
T2 (End) &
T2 (End) &
T1 &
T2 (End) &
T1 (End) &
First major escalation point: foundational, recovery, and performance-support categories reach Tier 2 by end gate, meaning resupply continuity, repeatable use-patterns, and phase-appropriate pre/intra/post support must be stable enough to survive travel, weather, and higher event consequence without creating volatility. 7.02 reaches Tier 1 as a legitimate clinician-reviewed support pathway when endocrine recovery, stress tolerance, or adaptation concerns warrant formal oversight. \\

07 &
Advanced Hybrid Overload (U.S. SOF Transformation Program) &
T2 &
T2 &
T2 (End) &
T2 &
T1 &
Hybrid density makes stimulant management and fatigue masking a real risk. 7.05 reaches Tier 2 by end gate only as a controlled-use system: standardized ceilings, cutoff rules, no rescue-dosing to salvage compliance, and no last-minute escalation during protected windows. 7.02 remains governed and clinician-mediated, not improvisational. \\

08 &
Structural Durability and Load Tolerance Phase &
T2 &
T2 &
T2 &
T2 &
T1 &
Durability phase shifts priority back toward tissue-support, sleep-support, hydration-support, and anti-friction consistency. Any cognitive or performance support that degrades sleep, hydration, appetite control, or foot-care discipline is operationally negative even if acute performance feels better. Stability matters more than aggression. \\

09 &
High-Volume Calisthenics and Stamina &
T2 &
T2 &
T2 &
T2 &
T1 &
High \textbf{CAL} volume rewards repeatable recovery and standardized energy management. GI stability, sleep protection, and hand/skin recovery matter more than aggressive ergogenic layering. Stack complexity should stay lower than the phase stress it is supposed to support. \\

10 &
Advanced Tactical Readiness &
T2 &
T2 &
T2 &
T2 &
T1 &
Multi-domain readiness requires a support stack that can survive mixed-event weeks without causing volatility. Phase rules reward boring repeatability: same products, same timing logic, same cutoff rules, same response to missed doses, and no “special event” additions close to gate windows. \\

11 &
Pre-Selection Conditioning and Recovery &
T2 &
T3 (End) &
T2 &
T2 &
T1 &
Recovery sensitivity is high enough that 7.06 reaches Tier 3 by end gate: sleep-protection, down-regulation, and recovery-support redundancy must be robust enough to protect final verification windows even when schedule stress rises. Energy support remains standardized, not aggressive, and must not bleed into sleep debt. \\

12 &
Final Pre-Selection Conditioning &
T3 (End) &
T3 &
T2 &
T2 &
T2 (End) &
Final pre-selection phase introduces Tier 3 foundational support and Tier 2 hormonal/adaptogen support by end gate where those systems are already established, clinically reviewed when appropriate, and stable under no-novelty rules. The requirement is reliability, tolerability, and repeatability under final verification pressure, not expansion for its own sake. \\

13 &
Full-Spectrum Readiness Consolidation &
T3 &
T3 &
T2 &
T2 &
T2 &
Consolidation phase holds the late-stage support stack steady. 7.03 and 7.06 remain Tier 3; 7.05 and 7.04 remain Tier 2; 7.02 remains Tier 2 only where already established, clinically governed as needed, and stable. No last-minute additions, removals, peaking changes, or “stronger” substitutions are admissible near final protected windows. \\

\end{longtblr}

\endgroup

\end{landscape}

\end{document}